\documentclass[a4paper,11pt]{article}

\usepackage[utf8]{inputenc}
\usepackage[T1]{fontenc}
\usepackage[english]{babel}
\usepackage{amsmath,amssymb,amsfonts}
\usepackage{braket}
\usepackage{geometry}
\geometry{margin=2.5cm}

\title{Proof of the Phase Kickback Phenomenon for the CNOT Gate}
\author{}
\date{}

\begin{document}

\maketitle

\section*{Introduction}

In this short document we show in detail how the circuit
\[
\mathrm{CNOT}\, (H \otimes HX)
\]
applied to the initial state \(\ket{00}\) implements the transformation
\[
\ket{+ -} \xrightarrow{\ \mathrm{CNOT}\ } \ket{- -},
\]
where the first qubit is the control and the second is the target.
This effect, in which the phase of the target is ``kicked back'' onto
the control, is known as \emph{phase kickback}.

\section*{Preliminaries}

We recall the definitions of the single-qubit \(\ket{\pm}\) states:
\begin{equation}
  \ket{+} = \frac{1}{\sqrt{2}}\bigl(\ket{0} + \ket{1}\bigr),\qquad
  \ket{-} = \frac{1}{\sqrt{2}}\bigl(\ket{0} - \ket{1}\bigr).
\end{equation}

The Hadamard gate \(H\) and Pauli \(X\) gate act on the computational basis as
\begin{align}
  H\ket{0} &= \ket{+}, & H\ket{1} &= \ket{-}, \\
  X\ket{0} &= \ket{1}, & X\ket{1} &= \ket{0}.
\end{align}

Moreover, \(\ket{-}\) is an eigenvector of \(X\) with eigenvalue \(-1\):
\begin{equation}
  X\ket{-}
  = X\left(\frac{\ket{0}-\ket{1}}{\sqrt{2}}\right)
  = \frac{\ket{1}-\ket{0}}{\sqrt{2}}
  = -\frac{\ket{0}-\ket{1}}{\sqrt{2}}
  = -\ket{-}.
  \label{eq:Xminus}
\end{equation}

The CNOT gate with the first qubit as control and the second as target acts as
\begin{equation}
  \mathrm{CNOT}\ket{00} = \ket{00},\quad
  \mathrm{CNOT}\ket{01} = \ket{01},\quad
  \mathrm{CNOT}\ket{10} = \ket{11},\quad
  \mathrm{CNOT}\ket{11} = \ket{10}.
\end{equation}

\section*{Preparing the state \texorpdfstring{$\ket{+ -}$}{|+ -⟩}}

We start from the initial state \(\ket{00}\) and apply the gates indicated in the circuit.

\subsection*{First qubit (control)}

On the first qubit we apply a Hadamard gate:
\begin{equation}
  (H \otimes I)\ket{00}
  = H\ket{0} \otimes \ket{0}
  = \ket{+}\ket{0}.
\end{equation}

\subsection*{Second qubit (target)}

On the second qubit we first apply \(X\) and then \(H\):
\begin{equation}
  (I \otimes X)\ket{0}
  = \ket{1},\qquad
  (I \otimes H X)\ket{0}
  = H X \ket{0}
  = H\ket{1}
  = \ket{-}.
\end{equation}

Overall, after the local gates we obtain the state
\begin{equation}
  \ket{\psi}
  = \ket{+}\ket{-}.
\end{equation}

Let us now write \(\ket{\psi}\) explicitly in the computational basis:
\begin{align}
  \ket{\psi}
  &= \ket{+}\ket{-} \\
  &= \left(\frac{\ket{0}+\ket{1}}{\sqrt{2}}\right)
     \otimes
     \left(\frac{\ket{0}-\ket{1}}{\sqrt{2}}\right) \\
  &= \frac{1}{2}\bigl(
        \ket{00} - \ket{01}
      + \ket{10} - \ket{11}
     \bigr).
  \label{eq:psi-explicit}
\end{align}

\section*{Action of the CNOT}

We now apply the CNOT gate to the state \(\ket{\psi}\) in
\eqref{eq:psi-explicit}:
\begin{align}
  \ket{\psi'}
  &= \mathrm{CNOT}\,\ket{\psi} \\
  &= \frac{1}{2}\bigl(
        \mathrm{CNOT}\ket{00}
      - \mathrm{CNOT}\ket{01}
      + \mathrm{CNOT}\ket{10}
      - \mathrm{CNOT}\ket{11}
     \bigr) \\
  &= \frac{1}{2}\bigl(
        \ket{00}
      - \ket{01}
      + \ket{11}
      - \ket{10}
     \bigr).
\end{align}

We group the terms according to the first qubit:
\begin{align}
  \ket{\psi'}
  &= \frac{1}{2}\Bigl[
        \ket{0}\bigl(\ket{0}-\ket{1}\bigr)
      + \ket{1}\bigl(\ket{1}-\ket{0}\bigr)
     \Bigr].
\end{align}

Note that
\begin{equation}
  \ket{0}-\ket{1} = \sqrt{2}\,\ket{-},\qquad
  \ket{1}-\ket{0} = -\sqrt{2}\,\ket{-},
\end{equation}
and therefore
\begin{align}
  \ket{\psi'}
  &= \frac{1}{2}\Bigl[
        \ket{0}\bigl(\sqrt{2}\ket{-}\bigr)
      + \ket{1}\bigl(-\sqrt{2}\ket{-}\bigr)
     \Bigr] \\
  &= \frac{\sqrt{2}}{2}
     \bigl(\ket{0} - \ket{1}\bigr)\ket{-} \\
  &= \frac{\sqrt{2}}{2}
     \bigl(\sqrt{2}\ket{-}\bigr)\ket{-} \\
  &= \ket{-}\ket{-}
   = \ket{- -}.
\end{align}

We have thus shown explicitly that
\begin{equation}
  \mathrm{CNOT}\,\ket{+ -} = \ket{- -}.
\end{equation}

\section*{Interpretation as phase kickback}

From the point of view of the second qubit, we see that it remains in the
state \(\ket{-}\) both before and after the CNOT.
However, from \eqref{eq:Xminus}, the state \(\ket{-}\) is an eigenstate of
\(X\) with eigenvalue \(-1\).
This means that when the control is \(\ket{1}\), the action of CNOT is
equivalent to multiplying the target state by a phase factor \(-1\).

Since the target is already in an eigenstate of \(X\), the net effect is not a
change of its state, but rather a \emph{shift} of this phase onto the control
qubit.
In other words, the \(\ket{1}\) component of the control acquires a minus sign
with respect to its \(\ket{0}\) component, turning \(\ket{+}\) into
\(\ket{-}\).
This transfer of phase from the target to the control is the
\emph{phase kickback} phenomenon.

\end{document}
